% Appendix section: PBPK-ABM algorithms and implementation listings
% Drop this into your dissertation with: % Appendix section: PBPK-ABM algorithms and implementation listings
% Drop this into your dissertation with: % Appendix section: PBPK-ABM algorithms and implementation listings
% Drop this into your dissertation with: % Appendix section: PBPK-ABM algorithms and implementation listings
% Drop this into your dissertation with: \input{docs/appendix_algorithms_listings}
% Assumes the preamble already includes algorithm2e, listings, and matlab-prettifier.

\section{Implementation Algorithms and Code Listings}
\label{sec:appendix-implementation-listings}

\subsection{Hybrid PBPK--ABM Workflow}
\begin{algorithm}[H]
\caption{Hybrid PBPK--ABM simulation loop with periodic coupling}
\label{alg:appendix-hybrid-workflow}
\DontPrintSemicolon
\KwIn{Dosing schedule, PBPK parameters, ABM initial tissue}
Initialize PBPK and ABM states\;
\For{$t \leftarrow 0$ \KwTo $t_{\max}$}{
  Advance substrate diffusion--decay on ABM mesh\;
  Advance all ABM cells (cycle, death, mechanics, motility)\;
  Interpolate PBPK concentrations at current $t$\;
  Inject $5$FU/FdUMP/FdUTP/FUTP fields into ABM microenvironment\;
  \If{$t \bmod 60\,\mathrm{min}=0$}{
    Record diagnostics and synchronized outputs\;
  }
  Update simulation clock\;
}
\end{algorithm}

\subsection{Ring-Based Tumour Seeding}
\begin{algorithm}[H]
\caption{Ring-based initialization of tumour and stromal/immune composition}
\label{alg:appendix-ring-seeding}
\DontPrintSemicolon
\KwIn{Ring table $\mathcal{R}$ with $(r_{\min},r_{\max},n_{type})$}
\ForEach{ring entry in $\mathcal{R}$}{
  \ForEach{cell type with count $n$ in this ring}{
    \For{$i\leftarrow 1$ \KwTo $n$}{
      Sample $\theta \sim U(0,2\pi)$\;
      Sample $r \sim \sqrt{U(r_{\min}^2,r_{\max}^2)}$ \tcp*{area-uniform annulus}
      Set $(x,y,z)=(r\cos\theta,r\sin\theta,0)$\;
      Place one agent of the selected cell type at $(x,y,z)$\;
    }
  }
}
Write all agents to \texttt{cellsX.csv}\;
\end{algorithm}

\subsection{Per-Cell Phenotype Update Under Drug Exposure}
\begin{algorithm}[H]
\caption{Per-cell drug response with S-phase sensitivity and apoptosis update}
\label{alg:appendix-phenotype-update}
\DontPrintSemicolon
\KwIn{Local $5$FU/FdUMP/FdUTP/FUTP, cycle phase, custom resistance data, $\Delta t$}
Read local metabolite concentrations at cell voxel\;
Apply ABC-efflux scaling to obtain effective intracellular exposure\;
\If{cell is in S-phase}{
  Compute DNA damage contributions from FdUMP and FdUTP\;
}
Compute RNA damage contribution from FUTP\;
Update total damage with repair term; clamp to $[0,1]$\;
\If{damage exceeds apoptosis threshold}{
  Increase apoptosis rate proportionally to excess damage\;
}
Update cumulative exposure and ABC-transporter induction\;
Apply updated death and resistance state to phenotype\;
\end{algorithm}

\subsection{Two-Mode Euler Stepping with Circadian DPD Modulation}
\begin{lstlisting}[style=Matlab-editor,caption={Fixed/adaptive Euler stepping with circadian Vmax scaling},label={lst:appendix-two-mode-euler}]
% Build timeline
if strcmp(params.solver_method,'fixed')
    time_min = (0:params.fixed_timestep_min:maxTime)';
else
    time_min = buildAdaptiveTimeline(dosingRegimen, ...
        params.adaptive_fine_timestep_min, ...
        params.adaptive_coarse_timestep_min, ...
        params.adaptive_window_before_min, ...
        params.adaptive_window_after_min, maxTime);
end

% Integrate ODE system
for t = 2:length(time_min)
    dt_actual = time_min(t) - time_min(t-1);
    currentTime = time_min(t-1);
    hourOfDay = mod(currentTime,1440)/60;
    DPD_factor = calculateCircadianDPD(hourOfDay, params);
    Vmax_current = params.Vmax_DPD * DPD_factor;

    dosingRate = calculateDosingRate(currentTime, dosingRegimen);
    dCdt = calculate5FUSystemODEs(C, t-1, dosingRate, ...
        DPD_factor, params, currentTime, dt_actual, logger);
    C = updateConcentrations(C, dCdt, t, dt_actual);  % Euler step
end
\end{lstlisting}

\subsection{PhysiCell Output Export and Feature Engineering}
\begin{algorithm}[H]
\caption{Pipeline from MultiCellDS snapshots to curve-fitter tables}
\label{alg:appendix-export-pipeline}
\DontPrintSemicolon
\KwIn{Set of PhysiCell \texttt{output*.xml} snapshots}
Parse each snapshot to recover cells, substrates, and metadata\;
Aggregate time-indexed summaries (counts, phases, densities)\;
Compute spatial/radial statistics by cell type\;
Compute microenvironment summary statistics\;
Assemble long-format and wide-format analysis tables\;
Generate curve-fitter table and highlight metrics\;
Write CSV outputs and metadata for reproducible post-processing\;
\end{algorithm}

\subsection{Monte Carlo Sensitivity Workflow}
\begin{algorithm}[H]
\caption{LHS-driven Monte Carlo sensitivity analysis for PBPK outputs}
\label{alg:appendix-mc-sensitivity}
\DontPrintSemicolon
\KwIn{$N$ samples, parameter ranges, dosing input}
Generate Latin Hypercube samples for selected PBPK parameters\;
Run one pre-flight simulation to validate pipeline\;
\ForEach{sample in parallel}{
  Apply parameter overrides\;
  Run PBPK simulation\;
  Extract AUC/Cmax and key endpoints\;
  Save chunk-level outputs and checkpoint state\;
}
Merge outputs; compute Spearman sensitivity metrics\;
Generate summary tables and figures\;
\end{algorithm}

\subsection{Numerical Audit and Verification}
\begin{lstlisting}[style=Matlab-editor,caption={Numerical audit checks for PBPK trajectory quality},label={lst:appendix-numerical-audit}]
results = run5FU_PBPK_Simulation(inputCsv, outputPrefix, paramOverrides);
t = results.time_min(:);
C = results.concentrations;

% 1) Timestep diagnostics
dt = diff(t);
report.time.min_dt = min(dt);
report.time.max_dt = max(dt);
report.time.n_nonpositive_dt = sum(dt <= 0);

% 2) Jump and clipping checks on key signals
for name = keySignals
    y = C.(name{1})(:);
    dy = diff(y);
    slope = dy ./ max(dt, eps);
    rel = abs(dy) ./ max(abs(y(1:end-1)), 1e-9);
    jumpMask = abs(slope) > jumpAbs | rel > jumpRel;
end

% 3) Cumulative monotonicity + mass-balance checks
report.mass.accounted_fraction = results.mass_balance.accounted_fraction;
\end{lstlisting}

\subsection{HPC Chunked Execution and Reproducibility}
\begin{lstlisting}[language=bash,caption={SLURM array submission pattern for chunked MC runs},label={lst:appendix-hpc-submission}]
TOTAL_SAMPLES=${TOTAL_SAMPLES:-10000}
CHUNK_SIZE=${CHUNK_SIZE:-100}
MAX_CONCURRENT=${MAX_CONCURRENT:-20}
ARRAY_TASKS=$(( (TOTAL_SAMPLES + CHUNK_SIZE - 1) / CHUNK_SIZE ))
ARRAY_SPEC="1-${ARRAY_TASKS}%${MAX_CONCURRENT}"

sbatch \
  --array="${ARRAY_SPEC}" \
  --export=ALL,CHUNK_SIZE=${CHUNK_SIZE},SEED_BASE=${SEED_BASE},OUT_ROOT=${OUT_ROOT} \
  scripts/hpc/run_MC_10k_array.sh

# Worker script executes one task/chunk and validates outputs
# before writing TASK_SUCCESS / TASK_FAILED markers.
\end{lstlisting}

\subsection{Unit-Test Evidence for Model Plumbing}
\begin{algorithm}[H]
\caption{Unit-test pattern for model configuration and substrate behavior checks}
\label{alg:appendix-unit-test-pattern}
\DontPrintSemicolon
\KwIn{Test fixture configuration, expected substrate/cell-definition behavior}
Set up minimal model state and input configuration\;
Execute targeted function or short simulation step\;
Assert expected changes in fields/definitions/outputs\;
Report pass/fail and write diagnostic traces if failed\;
Use focused cases for custom Dirichlet and cell-definition logic\;
\end{algorithm}

\paragraph{Appendix integration note.}
This section is intended as an implementation appendix: concise algorithmic summaries are paired with representative code fragments to preserve reproducibility while avoiding framework boilerplate.

% Assumes the preamble already includes algorithm2e, listings, and matlab-prettifier.

\section{Implementation Algorithms and Code Listings}
\label{sec:appendix-implementation-listings}

\subsection{Hybrid PBPK--ABM Workflow}
\begin{algorithm}[H]
\caption{Hybrid PBPK--ABM simulation loop with periodic coupling}
\label{alg:appendix-hybrid-workflow}
\DontPrintSemicolon
\KwIn{Dosing schedule, PBPK parameters, ABM initial tissue}
Initialize PBPK and ABM states\;
\For{$t \leftarrow 0$ \KwTo $t_{\max}$}{
  Advance substrate diffusion--decay on ABM mesh\;
  Advance all ABM cells (cycle, death, mechanics, motility)\;
  Interpolate PBPK concentrations at current $t$\;
  Inject $5$FU/FdUMP/FdUTP/FUTP fields into ABM microenvironment\;
  \If{$t \bmod 60\,\mathrm{min}=0$}{
    Record diagnostics and synchronized outputs\;
  }
  Update simulation clock\;
}
\end{algorithm}

\subsection{Ring-Based Tumour Seeding}
\begin{algorithm}[H]
\caption{Ring-based initialization of tumour and stromal/immune composition}
\label{alg:appendix-ring-seeding}
\DontPrintSemicolon
\KwIn{Ring table $\mathcal{R}$ with $(r_{\min},r_{\max},n_{type})$}
\ForEach{ring entry in $\mathcal{R}$}{
  \ForEach{cell type with count $n$ in this ring}{
    \For{$i\leftarrow 1$ \KwTo $n$}{
      Sample $\theta \sim U(0,2\pi)$\;
      Sample $r \sim \sqrt{U(r_{\min}^2,r_{\max}^2)}$ \tcp*{area-uniform annulus}
      Set $(x,y,z)=(r\cos\theta,r\sin\theta,0)$\;
      Place one agent of the selected cell type at $(x,y,z)$\;
    }
  }
}
Write all agents to \texttt{cellsX.csv}\;
\end{algorithm}

\subsection{Per-Cell Phenotype Update Under Drug Exposure}
\begin{algorithm}[H]
\caption{Per-cell drug response with S-phase sensitivity and apoptosis update}
\label{alg:appendix-phenotype-update}
\DontPrintSemicolon
\KwIn{Local $5$FU/FdUMP/FdUTP/FUTP, cycle phase, custom resistance data, $\Delta t$}
Read local metabolite concentrations at cell voxel\;
Apply ABC-efflux scaling to obtain effective intracellular exposure\;
\If{cell is in S-phase}{
  Compute DNA damage contributions from FdUMP and FdUTP\;
}
Compute RNA damage contribution from FUTP\;
Update total damage with repair term; clamp to $[0,1]$\;
\If{damage exceeds apoptosis threshold}{
  Increase apoptosis rate proportionally to excess damage\;
}
Update cumulative exposure and ABC-transporter induction\;
Apply updated death and resistance state to phenotype\;
\end{algorithm}

\subsection{Two-Mode Euler Stepping with Circadian DPD Modulation}
\begin{lstlisting}[style=Matlab-editor,caption={Fixed/adaptive Euler stepping with circadian Vmax scaling},label={lst:appendix-two-mode-euler}]
% Build timeline
if strcmp(params.solver_method,'fixed')
    time_min = (0:params.fixed_timestep_min:maxTime)';
else
    time_min = buildAdaptiveTimeline(dosingRegimen, ...
        params.adaptive_fine_timestep_min, ...
        params.adaptive_coarse_timestep_min, ...
        params.adaptive_window_before_min, ...
        params.adaptive_window_after_min, maxTime);
end

% Integrate ODE system
for t = 2:length(time_min)
    dt_actual = time_min(t) - time_min(t-1);
    currentTime = time_min(t-1);
    hourOfDay = mod(currentTime,1440)/60;
    DPD_factor = calculateCircadianDPD(hourOfDay, params);
    Vmax_current = params.Vmax_DPD * DPD_factor;

    dosingRate = calculateDosingRate(currentTime, dosingRegimen);
    dCdt = calculate5FUSystemODEs(C, t-1, dosingRate, ...
        DPD_factor, params, currentTime, dt_actual, logger);
    C = updateConcentrations(C, dCdt, t, dt_actual);  % Euler step
end
\end{lstlisting}

\subsection{PhysiCell Output Export and Feature Engineering}
\begin{algorithm}[H]
\caption{Pipeline from MultiCellDS snapshots to curve-fitter tables}
\label{alg:appendix-export-pipeline}
\DontPrintSemicolon
\KwIn{Set of PhysiCell \texttt{output*.xml} snapshots}
Parse each snapshot to recover cells, substrates, and metadata\;
Aggregate time-indexed summaries (counts, phases, densities)\;
Compute spatial/radial statistics by cell type\;
Compute microenvironment summary statistics\;
Assemble long-format and wide-format analysis tables\;
Generate curve-fitter table and highlight metrics\;
Write CSV outputs and metadata for reproducible post-processing\;
\end{algorithm}

\subsection{Monte Carlo Sensitivity Workflow}
\begin{algorithm}[H]
\caption{LHS-driven Monte Carlo sensitivity analysis for PBPK outputs}
\label{alg:appendix-mc-sensitivity}
\DontPrintSemicolon
\KwIn{$N$ samples, parameter ranges, dosing input}
Generate Latin Hypercube samples for selected PBPK parameters\;
Run one pre-flight simulation to validate pipeline\;
\ForEach{sample in parallel}{
  Apply parameter overrides\;
  Run PBPK simulation\;
  Extract AUC/Cmax and key endpoints\;
  Save chunk-level outputs and checkpoint state\;
}
Merge outputs; compute Spearman sensitivity metrics\;
Generate summary tables and figures\;
\end{algorithm}

\subsection{Numerical Audit and Verification}
\begin{lstlisting}[style=Matlab-editor,caption={Numerical audit checks for PBPK trajectory quality},label={lst:appendix-numerical-audit}]
results = run5FU_PBPK_Simulation(inputCsv, outputPrefix, paramOverrides);
t = results.time_min(:);
C = results.concentrations;

% 1) Timestep diagnostics
dt = diff(t);
report.time.min_dt = min(dt);
report.time.max_dt = max(dt);
report.time.n_nonpositive_dt = sum(dt <= 0);

% 2) Jump and clipping checks on key signals
for name = keySignals
    y = C.(name{1})(:);
    dy = diff(y);
    slope = dy ./ max(dt, eps);
    rel = abs(dy) ./ max(abs(y(1:end-1)), 1e-9);
    jumpMask = abs(slope) > jumpAbs | rel > jumpRel;
end

% 3) Cumulative monotonicity + mass-balance checks
report.mass.accounted_fraction = results.mass_balance.accounted_fraction;
\end{lstlisting}

\subsection{HPC Chunked Execution and Reproducibility}
\begin{lstlisting}[language=bash,caption={SLURM array submission pattern for chunked MC runs},label={lst:appendix-hpc-submission}]
TOTAL_SAMPLES=${TOTAL_SAMPLES:-10000}
CHUNK_SIZE=${CHUNK_SIZE:-100}
MAX_CONCURRENT=${MAX_CONCURRENT:-20}
ARRAY_TASKS=$(( (TOTAL_SAMPLES + CHUNK_SIZE - 1) / CHUNK_SIZE ))
ARRAY_SPEC="1-${ARRAY_TASKS}%${MAX_CONCURRENT}"

sbatch \
  --array="${ARRAY_SPEC}" \
  --export=ALL,CHUNK_SIZE=${CHUNK_SIZE},SEED_BASE=${SEED_BASE},OUT_ROOT=${OUT_ROOT} \
  scripts/hpc/run_MC_10k_array.sh

# Worker script executes one task/chunk and validates outputs
# before writing TASK_SUCCESS / TASK_FAILED markers.
\end{lstlisting}

\subsection{Unit-Test Evidence for Model Plumbing}
\begin{algorithm}[H]
\caption{Unit-test pattern for model configuration and substrate behavior checks}
\label{alg:appendix-unit-test-pattern}
\DontPrintSemicolon
\KwIn{Test fixture configuration, expected substrate/cell-definition behavior}
Set up minimal model state and input configuration\;
Execute targeted function or short simulation step\;
Assert expected changes in fields/definitions/outputs\;
Report pass/fail and write diagnostic traces if failed\;
Use focused cases for custom Dirichlet and cell-definition logic\;
\end{algorithm}

\paragraph{Appendix integration note.}
This section is intended as an implementation appendix: concise algorithmic summaries are paired with representative code fragments to preserve reproducibility while avoiding framework boilerplate.

% Assumes the preamble already includes algorithm2e, listings, and matlab-prettifier.

\section{Implementation Algorithms and Code Listings}
\label{sec:appendix-implementation-listings}

\subsection{Hybrid PBPK--ABM Workflow}
\begin{algorithm}[H]
\caption{Hybrid PBPK--ABM simulation loop with periodic coupling}
\label{alg:appendix-hybrid-workflow}
\DontPrintSemicolon
\KwIn{Dosing schedule, PBPK parameters, ABM initial tissue}
Initialize PBPK and ABM states\;
\For{$t \leftarrow 0$ \KwTo $t_{\max}$}{
  Advance substrate diffusion--decay on ABM mesh\;
  Advance all ABM cells (cycle, death, mechanics, motility)\;
  Interpolate PBPK concentrations at current $t$\;
  Inject $5$FU/FdUMP/FdUTP/FUTP fields into ABM microenvironment\;
  \If{$t \bmod 60\,\mathrm{min}=0$}{
    Record diagnostics and synchronized outputs\;
  }
  Update simulation clock\;
}
\end{algorithm}

\subsection{Ring-Based Tumour Seeding}
\begin{algorithm}[H]
\caption{Ring-based initialization of tumour and stromal/immune composition}
\label{alg:appendix-ring-seeding}
\DontPrintSemicolon
\KwIn{Ring table $\mathcal{R}$ with $(r_{\min},r_{\max},n_{type})$}
\ForEach{ring entry in $\mathcal{R}$}{
  \ForEach{cell type with count $n$ in this ring}{
    \For{$i\leftarrow 1$ \KwTo $n$}{
      Sample $\theta \sim U(0,2\pi)$\;
      Sample $r \sim \sqrt{U(r_{\min}^2,r_{\max}^2)}$ \tcp*{area-uniform annulus}
      Set $(x,y,z)=(r\cos\theta,r\sin\theta,0)$\;
      Place one agent of the selected cell type at $(x,y,z)$\;
    }
  }
}
Write all agents to \texttt{cellsX.csv}\;
\end{algorithm}

\subsection{Per-Cell Phenotype Update Under Drug Exposure}
\begin{algorithm}[H]
\caption{Per-cell drug response with S-phase sensitivity and apoptosis update}
\label{alg:appendix-phenotype-update}
\DontPrintSemicolon
\KwIn{Local $5$FU/FdUMP/FdUTP/FUTP, cycle phase, custom resistance data, $\Delta t$}
Read local metabolite concentrations at cell voxel\;
Apply ABC-efflux scaling to obtain effective intracellular exposure\;
\If{cell is in S-phase}{
  Compute DNA damage contributions from FdUMP and FdUTP\;
}
Compute RNA damage contribution from FUTP\;
Update total damage with repair term; clamp to $[0,1]$\;
\If{damage exceeds apoptosis threshold}{
  Increase apoptosis rate proportionally to excess damage\;
}
Update cumulative exposure and ABC-transporter induction\;
Apply updated death and resistance state to phenotype\;
\end{algorithm}

\subsection{Two-Mode Euler Stepping with Circadian DPD Modulation}
\begin{lstlisting}[style=Matlab-editor,caption={Fixed/adaptive Euler stepping with circadian Vmax scaling},label={lst:appendix-two-mode-euler}]
% Build timeline
if strcmp(params.solver_method,'fixed')
    time_min = (0:params.fixed_timestep_min:maxTime)';
else
    time_min = buildAdaptiveTimeline(dosingRegimen, ...
        params.adaptive_fine_timestep_min, ...
        params.adaptive_coarse_timestep_min, ...
        params.adaptive_window_before_min, ...
        params.adaptive_window_after_min, maxTime);
end

% Integrate ODE system
for t = 2:length(time_min)
    dt_actual = time_min(t) - time_min(t-1);
    currentTime = time_min(t-1);
    hourOfDay = mod(currentTime,1440)/60;
    DPD_factor = calculateCircadianDPD(hourOfDay, params);
    Vmax_current = params.Vmax_DPD * DPD_factor;

    dosingRate = calculateDosingRate(currentTime, dosingRegimen);
    dCdt = calculate5FUSystemODEs(C, t-1, dosingRate, ...
        DPD_factor, params, currentTime, dt_actual, logger);
    C = updateConcentrations(C, dCdt, t, dt_actual);  % Euler step
end
\end{lstlisting}

\subsection{PhysiCell Output Export and Feature Engineering}
\begin{algorithm}[H]
\caption{Pipeline from MultiCellDS snapshots to curve-fitter tables}
\label{alg:appendix-export-pipeline}
\DontPrintSemicolon
\KwIn{Set of PhysiCell \texttt{output*.xml} snapshots}
Parse each snapshot to recover cells, substrates, and metadata\;
Aggregate time-indexed summaries (counts, phases, densities)\;
Compute spatial/radial statistics by cell type\;
Compute microenvironment summary statistics\;
Assemble long-format and wide-format analysis tables\;
Generate curve-fitter table and highlight metrics\;
Write CSV outputs and metadata for reproducible post-processing\;
\end{algorithm}

\subsection{Monte Carlo Sensitivity Workflow}
\begin{algorithm}[H]
\caption{LHS-driven Monte Carlo sensitivity analysis for PBPK outputs}
\label{alg:appendix-mc-sensitivity}
\DontPrintSemicolon
\KwIn{$N$ samples, parameter ranges, dosing input}
Generate Latin Hypercube samples for selected PBPK parameters\;
Run one pre-flight simulation to validate pipeline\;
\ForEach{sample in parallel}{
  Apply parameter overrides\;
  Run PBPK simulation\;
  Extract AUC/Cmax and key endpoints\;
  Save chunk-level outputs and checkpoint state\;
}
Merge outputs; compute Spearman sensitivity metrics\;
Generate summary tables and figures\;
\end{algorithm}

\subsection{Numerical Audit and Verification}
\begin{lstlisting}[style=Matlab-editor,caption={Numerical audit checks for PBPK trajectory quality},label={lst:appendix-numerical-audit}]
results = run5FU_PBPK_Simulation(inputCsv, outputPrefix, paramOverrides);
t = results.time_min(:);
C = results.concentrations;

% 1) Timestep diagnostics
dt = diff(t);
report.time.min_dt = min(dt);
report.time.max_dt = max(dt);
report.time.n_nonpositive_dt = sum(dt <= 0);

% 2) Jump and clipping checks on key signals
for name = keySignals
    y = C.(name{1})(:);
    dy = diff(y);
    slope = dy ./ max(dt, eps);
    rel = abs(dy) ./ max(abs(y(1:end-1)), 1e-9);
    jumpMask = abs(slope) > jumpAbs | rel > jumpRel;
end

% 3) Cumulative monotonicity + mass-balance checks
report.mass.accounted_fraction = results.mass_balance.accounted_fraction;
\end{lstlisting}

\subsection{HPC Chunked Execution and Reproducibility}
\begin{lstlisting}[language=bash,caption={SLURM array submission pattern for chunked MC runs},label={lst:appendix-hpc-submission}]
TOTAL_SAMPLES=${TOTAL_SAMPLES:-10000}
CHUNK_SIZE=${CHUNK_SIZE:-100}
MAX_CONCURRENT=${MAX_CONCURRENT:-20}
ARRAY_TASKS=$(( (TOTAL_SAMPLES + CHUNK_SIZE - 1) / CHUNK_SIZE ))
ARRAY_SPEC="1-${ARRAY_TASKS}%${MAX_CONCURRENT}"

sbatch \
  --array="${ARRAY_SPEC}" \
  --export=ALL,CHUNK_SIZE=${CHUNK_SIZE},SEED_BASE=${SEED_BASE},OUT_ROOT=${OUT_ROOT} \
  scripts/hpc/run_MC_10k_array.sh

# Worker script executes one task/chunk and validates outputs
# before writing TASK_SUCCESS / TASK_FAILED markers.
\end{lstlisting}

\subsection{Unit-Test Evidence for Model Plumbing}
\begin{algorithm}[H]
\caption{Unit-test pattern for model configuration and substrate behavior checks}
\label{alg:appendix-unit-test-pattern}
\DontPrintSemicolon
\KwIn{Test fixture configuration, expected substrate/cell-definition behavior}
Set up minimal model state and input configuration\;
Execute targeted function or short simulation step\;
Assert expected changes in fields/definitions/outputs\;
Report pass/fail and write diagnostic traces if failed\;
Use focused cases for custom Dirichlet and cell-definition logic\;
\end{algorithm}

\paragraph{Appendix integration note.}
This section is intended as an implementation appendix: concise algorithmic summaries are paired with representative code fragments to preserve reproducibility while avoiding framework boilerplate.

% Assumes the preamble already includes algorithm2e, listings, and matlab-prettifier.

\section{Implementation Algorithms and Code Listings}
\label{sec:appendix-implementation-listings}

\subsection{Hybrid PBPK--ABM Workflow}
\begin{algorithm}[H]
\caption{Hybrid PBPK--ABM simulation loop with periodic coupling}
\label{alg:appendix-hybrid-workflow}
\DontPrintSemicolon
\KwIn{Dosing schedule, PBPK parameters, ABM initial tissue}
Initialize PBPK and ABM states\;
\For{$t \leftarrow 0$ \KwTo $t_{\max}$}{
  Advance substrate diffusion--decay on ABM mesh\;
  Advance all ABM cells (cycle, death, mechanics, motility)\;
  Interpolate PBPK concentrations at current $t$\;
  Inject $5$FU/FdUMP/FdUTP/FUTP fields into ABM microenvironment\;
  \If{$t \bmod 60\,\mathrm{min}=0$}{
    Record diagnostics and synchronized outputs\;
  }
  Update simulation clock\;
}
\end{algorithm}

\subsection{Ring-Based Tumour Seeding}
\begin{algorithm}[H]
\caption{Ring-based initialization of tumour and stromal/immune composition}
\label{alg:appendix-ring-seeding}
\DontPrintSemicolon
\KwIn{Ring table $\mathcal{R}$ with $(r_{\min},r_{\max},n_{type})$}
\ForEach{ring entry in $\mathcal{R}$}{
  \ForEach{cell type with count $n$ in this ring}{
    \For{$i\leftarrow 1$ \KwTo $n$}{
      Sample $\theta \sim U(0,2\pi)$\;
      Sample $r \sim \sqrt{U(r_{\min}^2,r_{\max}^2)}$ \tcp*{area-uniform annulus}
      Set $(x,y,z)=(r\cos\theta,r\sin\theta,0)$\;
      Place one agent of the selected cell type at $(x,y,z)$\;
    }
  }
}
Write all agents to \texttt{cellsX.csv}\;
\end{algorithm}

\subsection{Per-Cell Phenotype Update Under Drug Exposure}
\begin{algorithm}[H]
\caption{Per-cell drug response with S-phase sensitivity and apoptosis update}
\label{alg:appendix-phenotype-update}
\DontPrintSemicolon
\KwIn{Local $5$FU/FdUMP/FdUTP/FUTP, cycle phase, custom resistance data, $\Delta t$}
Read local metabolite concentrations at cell voxel\;
Apply ABC-efflux scaling to obtain effective intracellular exposure\;
\If{cell is in S-phase}{
  Compute DNA damage contributions from FdUMP and FdUTP\;
}
Compute RNA damage contribution from FUTP\;
Update total damage with repair term; clamp to $[0,1]$\;
\If{damage exceeds apoptosis threshold}{
  Increase apoptosis rate proportionally to excess damage\;
}
Update cumulative exposure and ABC-transporter induction\;
Apply updated death and resistance state to phenotype\;
\end{algorithm}

\subsection{Two-Mode Euler Stepping with Circadian DPD Modulation}
\begin{lstlisting}[style=Matlab-editor,caption={Fixed/adaptive Euler stepping with circadian Vmax scaling},label={lst:appendix-two-mode-euler}]
% Build timeline
if strcmp(params.solver_method,'fixed')
    time_min = (0:params.fixed_timestep_min:maxTime)';
else
    time_min = buildAdaptiveTimeline(dosingRegimen, ...
        params.adaptive_fine_timestep_min, ...
        params.adaptive_coarse_timestep_min, ...
        params.adaptive_window_before_min, ...
        params.adaptive_window_after_min, maxTime);
end

% Integrate ODE system
for t = 2:length(time_min)
    dt_actual = time_min(t) - time_min(t-1);
    currentTime = time_min(t-1);
    hourOfDay = mod(currentTime,1440)/60;
    DPD_factor = calculateCircadianDPD(hourOfDay, params);
    Vmax_current = params.Vmax_DPD * DPD_factor;

    dosingRate = calculateDosingRate(currentTime, dosingRegimen);
    dCdt = calculate5FUSystemODEs(C, t-1, dosingRate, ...
        DPD_factor, params, currentTime, dt_actual, logger);
    C = updateConcentrations(C, dCdt, t, dt_actual);  % Euler step
end
\end{lstlisting}

\subsection{PhysiCell Output Export and Feature Engineering}
\begin{algorithm}[H]
\caption{Pipeline from MultiCellDS snapshots to curve-fitter tables}
\label{alg:appendix-export-pipeline}
\DontPrintSemicolon
\KwIn{Set of PhysiCell \texttt{output*.xml} snapshots}
Parse each snapshot to recover cells, substrates, and metadata\;
Aggregate time-indexed summaries (counts, phases, densities)\;
Compute spatial/radial statistics by cell type\;
Compute microenvironment summary statistics\;
Assemble long-format and wide-format analysis tables\;
Generate curve-fitter table and highlight metrics\;
Write CSV outputs and metadata for reproducible post-processing\;
\end{algorithm}

\subsection{Monte Carlo Sensitivity Workflow}
\begin{algorithm}[H]
\caption{LHS-driven Monte Carlo sensitivity analysis for PBPK outputs}
\label{alg:appendix-mc-sensitivity}
\DontPrintSemicolon
\KwIn{$N$ samples, parameter ranges, dosing input}
Generate Latin Hypercube samples for selected PBPK parameters\;
Run one pre-flight simulation to validate pipeline\;
\ForEach{sample in parallel}{
  Apply parameter overrides\;
  Run PBPK simulation\;
  Extract AUC/Cmax and key endpoints\;
  Save chunk-level outputs and checkpoint state\;
}
Merge outputs; compute Spearman sensitivity metrics\;
Generate summary tables and figures\;
\end{algorithm}

\subsection{Numerical Audit and Verification}
\begin{lstlisting}[style=Matlab-editor,caption={Numerical audit checks for PBPK trajectory quality},label={lst:appendix-numerical-audit}]
results = run5FU_PBPK_Simulation(inputCsv, outputPrefix, paramOverrides);
t = results.time_min(:);
C = results.concentrations;

% 1) Timestep diagnostics
dt = diff(t);
report.time.min_dt = min(dt);
report.time.max_dt = max(dt);
report.time.n_nonpositive_dt = sum(dt <= 0);

% 2) Jump and clipping checks on key signals
for name = keySignals
    y = C.(name{1})(:);
    dy = diff(y);
    slope = dy ./ max(dt, eps);
    rel = abs(dy) ./ max(abs(y(1:end-1)), 1e-9);
    jumpMask = abs(slope) > jumpAbs | rel > jumpRel;
end

% 3) Cumulative monotonicity + mass-balance checks
report.mass.accounted_fraction = results.mass_balance.accounted_fraction;
\end{lstlisting}

\subsection{HPC Chunked Execution and Reproducibility}
\begin{lstlisting}[language=bash,caption={SLURM array submission pattern for chunked MC runs},label={lst:appendix-hpc-submission}]
TOTAL_SAMPLES=${TOTAL_SAMPLES:-10000}
CHUNK_SIZE=${CHUNK_SIZE:-100}
MAX_CONCURRENT=${MAX_CONCURRENT:-20}
ARRAY_TASKS=$(( (TOTAL_SAMPLES + CHUNK_SIZE - 1) / CHUNK_SIZE ))
ARRAY_SPEC="1-${ARRAY_TASKS}%${MAX_CONCURRENT}"

sbatch \
  --array="${ARRAY_SPEC}" \
  --export=ALL,CHUNK_SIZE=${CHUNK_SIZE},SEED_BASE=${SEED_BASE},OUT_ROOT=${OUT_ROOT} \
  scripts/hpc/run_MC_10k_array.sh

# Worker script executes one task/chunk and validates outputs
# before writing TASK_SUCCESS / TASK_FAILED markers.
\end{lstlisting}

\subsection{Unit-Test Evidence for Model Plumbing}
\begin{algorithm}[H]
\caption{Unit-test pattern for model configuration and substrate behavior checks}
\label{alg:appendix-unit-test-pattern}
\DontPrintSemicolon
\KwIn{Test fixture configuration, expected substrate/cell-definition behavior}
Set up minimal model state and input configuration\;
Execute targeted function or short simulation step\;
Assert expected changes in fields/definitions/outputs\;
Report pass/fail and write diagnostic traces if failed\;
Use focused cases for custom Dirichlet and cell-definition logic\;
\end{algorithm}

\paragraph{Appendix integration note.}
This section is intended as an implementation appendix: concise algorithmic summaries are paired with representative code fragments to preserve reproducibility while avoiding framework boilerplate.
